\chapter{Revisit 3}

\begin{center}
    {\Huge \bfseries Andi - Japanese Handbook}\\[0.5em]
    {\LARGE \color{boxframe} \textbf{Bab Spesial 3: Keahlian Bertahan Hidup (Advanced)}}
\end{center}
\vspace{1em}

% ================= THEME 1 =================
\begin{lessonbox}{1. Menerima, Memberi \& Hadiah (Giving \& Receiving 1-5)}
    \textbf{Topik: Giving and receiving \#1 - \#5} \\
    Di Jepang, budaya memberi oleh-oleh (\textit{Omiyage}) sangat kental! 
    \begin{itemize}
        \item \textbf{あげる (Ageru)} = Saya memberi ke orang lain.
        \item \textbf{もらう (Morau)} = Saya menerima dari orang lain.
        \item \textbf{くれる (Kureru)} = Orang lain memberi ke saya!
    \end{itemize}
    \textbf{Tips Turis:} Kalau ada orang Jepang yang sangat baik membantumu membawakan koper atau memberimu peta, kamu bisa bilang \textit{"Tetsudatte kurete, arigatou!"} (Terima kasih sudah memberiku bantuan!). Jika kamu ingin memberikan suvenir dari Indonesia, katakan \textit{"Kore, agemasu!"} (Ini, saya berikan untukmu!).
\end{lessonbox}

\begin{convobox}{1}
    \noindent
    \jpword{ともだち}{友達}{Tomodachi} \jpkana{に}{ni} 
    \jpword{おみやげ}{お土産}{o-miyage} \jpkana{を}{o} 
    \jpkana{あげます。}{agemasu.}

    \vspace{1.5em}
    \noindent{\Large \color{articolor} \textbf{Arti:} Saya memberikan oleh-oleh kepada teman.}
\end{convobox}

\begin{convobox}{2}
    \noindent
    \jpword{てつだ}{手伝}{Tetsuda}\jpkana{って}{tte} 
    \jpkana{くれて、}{kurete,} 
    \jpkana{ありがとう。}{arigatou.}

    \vspace{1.5em}
    \noindent{\Large \color{articolor} \textbf{Arti:} Terima kasih sudah (memberikan) bantuan kepadaku.}
\end{convobox}

% ================= THEME 2 =================
\begin{lessonbox}{2. Menjelaskan Benda secara Detail (Sifat \& Klausa)}
    \textbf{Topik: More の usages, Relative Clauses, Describing nouns/actions, こと+が} \\
    Bagaimana cara bilang "Kereta yang pergi ke Tokyo" atau "Sushi yang dimakan Ayah"? Di Jepang, \textbf{semua penjelasannya ditaruh di DEPAN kata benda}! 
    \begin{itemize}
        \item \textit{Toukyou ni iku} (Pergi ke Tokyo) + \textit{Densha} (Kereta) = \textbf{Toukyou ni iku densha} (Kereta yang pergi ke Tokyo).
    \end{itemize}
    \textbf{Tips Turis:} Kalau kamu kehilangan barang, kamu bisa menjelaskannya dengan detail. Misalnya mencari tas: \textit{"Watashi ga katta kaban"} (Tas yang saya beli). Oh ya, partikel \textbf{の (No)} juga bisa menggantikan benda agar kita tidak repot mengulang kata. "Saya suka yang merah" = \textit{"Akai \textbf{no} ga suki desu"}.
\end{lessonbox}

\begin{convobox}{3}
    \noindent
    \jpword{とうきょう}{東京}{Toukyou} \jpkana{に}{ni} 
    \jpword{い}{行}{i}\jpkana{く}{ku} 
    \jpword{でんしゃ}{電車}{densha} \jpkana{は}{wa} 
    \jpkana{どれ}{dore} \jpkana{ですか？}{desu ka?}

    \vspace{1.5em}
    \noindent{\Large \color{articolor} \textbf{Arti:} Kereta yang pergi ke Tokyo yang mana?}
\end{convobox}

\begin{convobox}{4}
    \noindent
    \jpword{あか}{赤}{Aka}\jpkana{い}{i} \jpkana{の}{no} \jpkana{を}{o} 
    \jpword{くだ}{下}{kuda}\jpkana{さい。}{sai.}

    \vspace{1.5em}
    \noindent{\Large \color{articolor} \textbf{Arti:} Tolong berikan (barang) YANG merah.}
\end{convobox}

% ================= THEME 3 =================
\begin{lessonbox}{3. Bisa, Boleh, dan Harus! (Aturan \& Kemampuan)}
    \textbf{Topik: Potential form, Must \& Must Not, Commands, Suggestions, Explanations (んです)} \\
    Saat jalan-jalan, sangat penting memahami papan peringatan dan aturan!
    \begin{itemize}
        \item \textbf{Bisa (Potential):} Ubah \textit{~masu} menjadi \textit{~emasu}. Contoh: \textit{Taberaremasu} (Bisa dimakan).
        \item \textbf{Harus (Must):} \textit{〜なければならない (〜nakereba naranai)}. "Saya harus pergi!"
        \item \textbf{Saran (Suggestions):} \textit{〜ほうがいい (〜hou ga ii)}. "Lebih baik bawa payung!"
    \end{itemize}
    \textbf{Tips Turis:} Kalau kamu tersesat dan butuh menjelaskan situasimu, gunakan \textbf{〜んです (〜n desu)} di akhir kalimatmu. Ini memberi kesan "Tolong dengarkan alasanku!". Contoh: \textit{"Michi ni mayotta n desu!"} (Aku tersesat!).
\end{lessonbox}

\begin{convobox}{5}
    \noindent
    \jpkana{ここで}{Koko de} 
    \jpword{しゃしん}{写真}{shashin} \jpkana{を}{o} 
    \jpword{と}{撮}{to}\jpkana{ってはいけません。}{tte wa ikemasen.}

    \vspace{1.5em}
    \noindent{\Large \color{articolor} \textbf{Arti:} Tidak boleh (Dilarang) mengambil foto di sini.}
\end{convobox}

\begin{convobox}{6}
    \noindent
    \jpword{きっぷ}{切符}{Kippu} \jpkana{を}{o} 
    \jpword{か}{買}{ka}\jpkana{わなければなりません。}{wanakereba narimasen.}

    \vspace{1.5em}
    \noindent{\Large \color{articolor} \textbf{Arti:} (Saya) harus membeli tiket kereta.}
\end{convobox}

% ================= THEME 4 =================
\begin{lessonbox}{4. Bahasa Dewa ala Jepang (Honorific \& Humble)}
    \textbf{Topik: Honorific speech, Humble speech, Passive voice, Causative verbs} \\
    Pernah dengar pengumuman stasiun atau kasir yang bahasanya terdengar asing dan panjang sekali? Itu adalah \textbf{敬語 (Keigo / Bahasa Hormat)}!
    \begin{itemize}
        \item \textbf{Honorific (Meninggikan tamu):} \textit{O-machi kudasai} (Tolong tunggu).
        \item \textbf{Humble (Merendahkan diri):} \textit{Itashimasu} (Saya akan melakukan).
    \end{itemize}
    \textbf{Tips Turis:} Kamu \textbf{TIDAK PERLU} bisa berbicara dengan Keigo! Cukup ucapkan bentuk \textit{〜masu} biasa. Namun, biasakan telingamu mendengarnya agar kamu paham maksud pelayan toko. Oh ya, kalau kakimu tidak sengaja terinjak di kereta, kamu bisa pakai \textbf{Pasif (Passive Voice)}: \textit{"Ashi o fUmaremashita!"} (Kakiku terinjak!).
\end{lessonbox}

\begin{convobox}{7}
    \noindent
    \jpkana{しょうしょう}{Shoushou} \jpkana{お}{o}\jpword{ま}{待}{ma}\jpkana{ち}{chi} 
    \jpword{くだ}{下}{kuda}\jpkana{さい。}{sai.}

    \vspace{1.5em}
    \noindent{\Large \color{articolor} \textbf{Arti:} Mohon tunggu sebentar. (Bahasa Honorific Kasir)}
\end{convobox}

\begin{convobox}{8}
    \noindent
    \jpword{わたし}{私}{Watashi} \jpkana{が}{ga} 
    \jpword{あんない}{案内}{annai} \jpkana{いたします。}{itashimasu.}

    \vspace{1.5em}
    \noindent{\Large \color{articolor} \textbf{Arti:} Saya yang akan memandu Anda. (Bahasa Humble Staf)}
\end{convobox}

% ================= THEME 5 =================
\begin{lessonbox}{5. Waktu, Pergerakan, dan Batas (Time \& Movement)}
    \textbf{Topik: Going \& Coming, Noun Suffixes (まで, Time, でも, Specification, Other)} \\
    Bagaimana cara bilang pergerakan dari A ke B? 
    \begin{itemize}
        \item \textbf{〜まで (〜made)} = Sampai. Contoh: \textit{Toukyou made} (Sampai Tokyo).
        \item \textbf{〜ていく / 〜てくる (〜te iku / 〜te kuru)} = Pergi / Datang. \textit{Katte kimasu!} (Aku beli dulu ya lalu kembali ke sini!).
        \item \textbf{〜でも (〜demo)} = "Atau semacamnya". \textit{Koohii demo nomimasu ka?} (Mau minum kopi atau semacamnya?).
    \end{itemize}
    \textbf{Tips Turis:} Gunakan \textit{〜made} saat bertanya pada supir taksi: \textit{"Kono hoteru \textbf{made} onegaishimasu"} (Tolong antarkan \textbf{sampai} hotel ini).
\end{lessonbox}

\begin{convobox}{9}
    \noindent
    \jpword{えき}{駅}{Eki} \jpkana{まで}{made} 
    \jpword{い}{行}{i}\jpkana{ってきます。}{tte kimasu.}

    \vspace{1.5em}
    \noindent{\Large \color{articolor} \textbf{Arti:} Saya akan pergi ke stasiun (lalu akan kembali lagi ke sini).}
\end{convobox}

% ================= THEME 6 =================
\begin{lessonbox}{6. Menyambung Kalimat \& Alasan (Conjunctions)}
    \textbf{Topik: ため/ために, Cause and Effect, のに/だけど, から, Linking sentences, Conjunctions 1-2} \\
    Agar kamu bisa bercerita layaknya penduduk lokal, gunakan kata sambung (Conjunctions)!
    \begin{itemize}
        \item \textbf{〜ために (〜tame ni)} = Demi / Untuk. \textit{Kazoku no tame ni} (Demi keluarga).
        \item \textbf{だから (Dakara) / から (Kara)} = Oleh karena itu / Karena.
        \item \textbf{なのに (Na noni) / だけど (Dakedo)} = Padahal / Tetapi.
    \end{itemize}
    \textbf{Tips Turis:} \textit{Dakedo} sangat sering dipakai orang Jepang di bahasa sehari-hari. "Sepatu ini mahal, \textit{dakedo} nyaman!"
\end{lessonbox}

\begin{convobox}{10}
    \noindent
    \jpword{あめ}{雨}{Ame} \jpkana{が}{ga} 
    \jpword{ふ}{降}{fu}\jpkana{っている}{tte iru} 
    \jpkana{から、}{kara,} 
    \jpkana{タクシーで}{takushii de} 
    \jpword{い}{行}{i}\jpkana{きましょう。}{kimashou.}

    \vspace{1.5em}
    \noindent{\Large \color{articolor} \textbf{Arti:} KARENA sedang turun hujan, mari kita pergi naik taksi.}
\end{convobox}

% ================= THEME 7 =================
\begin{lessonbox}{7. Kelihatannya \& Katanya (Senses \& Guesses)}
    \textbf{Topik: Seems らしい/よう/そう, Hearsay, Expectations, Intentions, Supposition} \\
    Melihat kue stroberi di etalase kafe? Kamu bisa bilang \textbf{「美味しそう！」 (Oishisou! - Kelihatannya enak!)}. 
    \begin{itemize}
        \item \textbf{〜そう (〜sou)} = Kelihatannya... (Dari apa yang kita lihat).
        \item \textbf{〜らしい (〜rashii) / 〜そうです (〜sou desu)} = Katanya / Kudengar... (Dari gosip atau berita).
        \item \textbf{〜つもり (〜tsumori)} = Rencana / Niat.
    \end{itemize}
    \textbf{Tips Turis:} Kalau besok ada rencana ke Gunung Fuji, kamu bisa bilang \textit{"Ashita, Fuji-san ni iku \textbf{tsumori} desu!"} (Besok, rencananya mau pergi ke Gunung Fuji!).
\end{lessonbox}

\begin{convobox}{11}
    \noindent
    \jpkana{この}{Kono} \jpkana{ケーキ}{keeki}\jpkana{、}{,} 
    \jpkana{おいし}{oishi}\jpkana{そうですね！}{sou desu ne!}

    \vspace{1.5em}
    \noindent{\Large \color{articolor} \textbf{Arti:} Kue ini, kelihatan enak ya!}
\end{convobox}

\begin{convobox}{12}
    \noindent
    \jpword{あした}{明日}{Ashita} \jpkana{は}{wa} 
    \jpword{ゆき}{雪}{yuki} \jpkana{が}{ga} 
    \jpword{ふ}{降}{fu}\jpkana{る}{ru} 
    \jpkana{そうです。}{sou desu.}

    \vspace{1.5em}
    \noindent{\Large \color{articolor} \textbf{Arti:} KATANYA besok akan turun salju (Dengar dari berita cuaca).}
\end{convobox}

% ================= THEME 8 =================
\begin{lessonbox}{8. Pengandaian \& Keputusan (Conditionals)}
    \textbf{Topik: Conditional たら/ば, ことにする/なる, ように} \\
    Bagaimana jika hujan? Bagaimana jika tiketnya habis? Gunakan \textbf{〜たら (〜tara)} atau \textbf{〜ば (〜ba)} untuk "Jika/Kalau".
    \begin{itemize}
        \item \textbf{〜たら (〜tara)} = Jika (Paling sering dipakai dan paling mudah!). \textit{Ame ga futtara...} (Kalau hujan...).
        \item \textbf{〜ことにする (〜koto ni suru)} = Saya memutuskan untuk... \textit{Kore o kau koto ni shimasu!} (Saya putuskan beli yang ini!).
    \end{itemize}
\end{lessonbox}

\begin{convobox}{13}
    \noindent
    \jpword{にほん}{日本}{Nihon} \jpkana{に}{ni} 
    \jpword{つ}{着}{tsu}\jpkana{いたら、}{itara,} 
    \jpkana{ラーメン}{raamen} \jpkana{を}{o} 
    \jpword{た}{食}{ta}\jpkana{べましょう！}{bemashou!}

    \vspace{1.5em}
    \noindent{\Large \color{articolor} \textbf{Arti:} JIKA (nanti) sudah tiba di Jepang, mari kita makan ramen!}
\end{convobox}

% ================= THEME 9 =================
\begin{lessonbox}{9. Memulai, Selesai, \& Kondisi (Actions in Progress)}
    \textbf{Topik: ところ, Completing (てしまう), Starting (はじめる), Verb States, Volitional Expanded, Advanced Questions, Comparisons 2} \\
    Terakhir, mari mengekspresikan tindakanmu dengan lebih tajam!
    \begin{itemize}
        \item \textbf{〜ところ (〜tokoro)} = Baru saja mau / Sedang / Baru saja selesai. \textit{Ima, iku tokoro desu!} (Sekarang, aku baru saja mau berangkat!).
        \item \textbf{〜てしまう (〜te shimau)} = Selesai (tuntas), atau "Waduh/Gawat!" (Tidak sengaja).
        \item \textbf{〜はじめる (〜hajimeru)} = Mulai melakukan sesuatu. \textit{Ame ga furi-hajimemashita} (Hujan mulai turun).
    \end{itemize}
    \textbf{Tips Turis Darurat:} Kalau kamu tidak sengaja menjatuhkan atau menghilangkan paspor/tiket, gunakan \textit{te shimaimashita} untuk menunjukkan rasa menyesal/gawat. \textit{"Kippu o nakushite shimaimashita!"} (Waduh, tiket saya hilang!). Petugas akan sangat mengerti kepanikanmu!
\end{lessonbox}

\begin{convobox}{14}
    \noindent
    \jpword{きっぷ}{切符}{Kippu} \jpkana{を}{o} 
    \jpword{な}{失}{na}\jpkana{くして}{kushite} 
    \jpkana{しまいました。}{shimaimashita.}

    \vspace{1.5em}
    \noindent{\Large \color{articolor} \textbf{Arti:} Waduh, saya (tidak sengaja) menghilangkan tiket saya.}
\end{convobox}

\begin{convobox}{15}
    \noindent
    \jpword{いま}{今}{Ima}\jpkana{、}{,} 
    \jpword{ひこうき}{飛行機}{hikouki} \jpkana{に}{ni} 
    \jpword{の}{乗}{no}\jpkana{る}{ru} 
    \jpkana{ところです！}{tokoro desu!}

    \vspace{1.5em}
    \noindent{\Large \color{articolor} \textbf{Arti:} Sekarang, saya baru saja mau naik ke pesawat!}
\end{convobox}

% --- SUMMARY BOX CLOSING ---
\vspace{2em}
\begin{tcolorbox}[
    colback=red!5!white,
    colframe=red!60!black,
    boxrule=3pt,
    arc=5mm,
    title={\huge\bfseries \sffamily 🏆 LULUS TINGKAT MASTER! 🏆},
    fonttitle=\color{white},
    attach boxed title to top center={yshift=-4mm},
    boxed title style={colback=red!60!black, rounded corners},
    left=5mm, right=5mm, top=7mm, bottom=5mm
]
\Large \textbf{Luar Biasa!} Ayah, Bunda, dan Adik-adik kini sudah memegang "Sabuk Hitam" untuk bertahan hidup di Jepang! Ke-50 topik tingkat lanjut ini memang menantang, tapi dengan mengingat tips-tips kuncinya:
\begin{itemize}
    \item Dengarkan kata kuncinya (Jangan panik jika mendengar \textit{Keigo} / bahasa dewa kasir).
    \item Gunakan \textbf{〜んです (n desu)} jika butuh menjelaskan alasan darurat.
    \item Gunakan \textbf{〜てしまいました (te shimaimashita)} jika tidak sengaja melakukan kesalahan.
\end{itemize}
Sekarang, liburan kalian di Jepang pasti akan penuh dengan petualangan bahasa yang luar biasa. \textbf{ご旅行、楽しんでくださいね！ (Go-ryokou, tanoshinde kudasai ne - Nikmati liburannya ya!)}
\end{tcolorbox}

\newpage